\documentclass[12pt]{article}
\usepackage[a4paper,margin=2.5cm]{geometry}
\usepackage{times}
\usepackage{hyperref}
\usepackage{parskip}
\pagestyle{empty}

\begin{document}

\begin{flushright}
Chengdu, \today
\end{flushright}

\noindent
Editorial Office\\
\emph{Journal of Power Sources}\\
Elsevier

\bigskip
\noindent Dear Editors,

Please consider the enclosed manuscript entitled \textbf{``A Modified Shepherd Electrochemical Model with Coupled Temperature--Aging Dynamics for Lithium-Ion Cells under Constant-Current and Random-Walk Discharge''} for publication in the \emph{Journal of Power Sources} as a research article.

The work develops a modified Shepherd-based electrochemical modelling framework for lithium-ion cells under the dynamic, heterogeneous discharge regimes characteristic of portable electronic systems, and validates it on two independent public datasets.
Three contributions are central to the manuscript and align directly with the scope of \emph{Journal of Power Sources}:
\begin{enumerate}
    \item \textbf{An improved Shepherd voltage formulation with retained physical interpretability.}
    The five-parameter modified Shepherd equation preserves the one-to-one mapping between coefficients and electrochemical phenomena (open-circuit potential, mass-transport polarisation, ohmic loss, exponential transient zone), and is coupled with a logistic capacity--temperature submodel, an Arrhenius resistance--temperature submodel, and power-law cycle/$\sqrt{t}$ calendar aging into a single continuous-time differential-algebraic system.
    On the XJTU lithium-ion dataset the formulation achieves a voltage-tracking RMSE of 17.85\,mV under constant-current discharge, reducing error by over 55\% relative to a Th\'{e}venin-1RC ECM and 75\% relative to a Nernst baseline.
    \item \textbf{Cross-batch generalisation across C-rates.}
    Parameters fitted on XJTU Batch-1 (2\,C, 15 cells) are applied directly to Batch-2 (3\,C) without re-identification, with a statistically non-significant RMSE gap (95\% bootstrap CI spanning zero).
    To our knowledge this transferability test---explicitly requested in many \emph{Journal of Power Sources} reviews of empirical voltage models---has not previously been reported for Shepherd-type formulations on lithium-ion chemistries.
    \item \textbf{Dynamic-load validation on the NASA Randomized Battery Usage Dataset.}
    Using the same parameters identified from constant-current data, the framework is validated against random-walk currents spanning 0.5--4\,A over 200 cycles, yielding a voltage RMSE of 63.5\,mV with no re-fitting.
    This demonstrates that the Shepherd formulation extends from idealised CC discharges to realistic time-varying loads---a key concern for any battery model intended for deployment in modern portable systems.
\end{enumerate}

A systematic ablation across a $3{\times}3$ temperature--aging grid further quantifies the non-additive coupling between submodels (temperature contributing up to 70\% TTE error, aging up to 55\%, with non-trivial interaction terms).
A receding-horizon MPC strategy built on the Shepherd prediction core illustrates one downstream application, extending simulated runtime by up to 58\% with a computational cost of 0.009\,ms per simulated second, suitable for embedded BMS deployment.
The modelling and validation results stand independently of this control study.

We believe the manuscript fits the scope of \emph{Journal of Power Sources} because it advances the empirical/electrochemical modelling of lithium-ion cells, provides quantitative cross-batch and dynamic-load validation on widely used public datasets (XJTU, NASA), and reports parameters and code that the lithium-ion battery community can directly reuse.

This manuscript has not been published previously and is not under consideration elsewhere.
All authors have approved the manuscript and agree with its submission.
No conflicts of interest exist.

\bigskip
\noindent\textbf{Suggested reviewers}
\begin{itemize}
    \item Prof.\ David A.\ Howey, Department of Engineering Science, University of Oxford, UK --- expertise in lithium-ion modelling, ECM identification, and Gaussian-process battery state estimation; cited in the manuscript (Birkl, Richardson).
    \item Prof.\ Dirk Uwe Sauer, ISEA, RWTH Aachen University, Germany --- expertise in lithium-ion aging modelling and characterisation; corresponding author of the holistic NMC aging model cited in the manuscript (Schmalstieg \emph{et al.}, JPS 2014).
    \item Prof.\ Gregory L.\ Plett, Department of Electrical and Computer Engineering, University of Colorado Colorado Springs, USA --- expertise in ECM-based battery state estimation for HEV battery packs; cited in the manuscript (Plett, JPS 2004).
    \item Prof.\ Xiaosong Hu, College of Mechanical and Vehicle Engineering, Chongqing University, China --- expertise in comparative ECM studies and battery management; cited in the manuscript (Hu \emph{et al.}, JPS 2012).
    \item Prof.\ Rui Xiong, School of Mechanical Engineering, Beijing Institute of Technology, China --- expertise in lithium-ion state-of-health monitoring and battery management systems; cited in the manuscript (Xiong \emph{et al.}, JPS 2018).
\end{itemize}
We have no conflicting relationship with any of the suggested reviewers.

\bigskip
\noindent Cordially yours,

\medskip
\noindent Prof.\ Zhiqiang Li\\
Corresponding author

\bigskip
\noindent\textit{P.S.}\ For any questions and future correspondence, the corresponding and contact author will be:\\
Prof.\ Zhiqiang Li\\
E-mail: \texttt{zhiqiangli@scu.edu.cn}\\
College of Physics, Sichuan University\\
No.\ 24 South Section 1, Yihuan Road\\
Chengdu 610064, China

\end{document}
